\ExerciseSection

\begin{enumerate}
\def\labelenumi{\arabic{enumi})}
\tightlist
\item
  伪度量 \(f\) （在群 \(G\) 上，乘法记号）称为左不变的（相应地，右不变的），如果 \(f(zx, zy) = f(x, y)\) {[}相应地，\(f(xz, yz) = f(x, y)\) {]} 对所有 \(x, y, z\) ∈ \(G\) 成立。
\end{enumerate}

\begin{enumerate}
\def\labelenumi{\alph{enumi})}
\tightlist
\item
  若 \(f\) 是一个左不变伪度量，则实值函数 \(g(x) = f(e, x)\) （在 \(\mathbf{G}\) 上；\(e\) 是 \(\mathbf{G}\) 的单位元）满足下列条件：
\end{enumerate}

\begin{enumerate}
\def\labelenumi{(\roman{enumi})}
\item
  \(g(x)\geqslant \mathrm{0}\) 对所有 \(x\in \mathbf{G}\) 成立，且 \(g(e) = \mathrm{0};\)
\item
  \(g(x^{-1}) = g(x)\) ；
\item
  \(g(xy) \leqslant g(x) + g(y)\) 。
\end{enumerate}

反之，若 \(g\) 是 \(G\) 上满足这些条件的任一实值函数，则 \(f(x, y) = g(x^{-1}y)\) 是 \(G\) 上的一个左不变伪度量。\(b)\) 拓扑 \(\mathfrak{G}\) ——由左不变伪度量组成的一个饱和族 \((f_t)\) （在群 \(G\) 上）所定义——与 \(G\) 的群结构相容，当且仅当对每个 \(a \in G\) 、每个指标 \(t\) 与每个实数 \(\alpha > 0\) ，存在一个指标 \(x\) 与一个实数 \(\beta > 0\) ，使得关系 \(f_x(e, x) \leqslant \beta\) 蕴含 \(f_t(e, axa^{-1}) \leqslant \alpha\) 。若此条件满足，则 \(G\) 上由族 \(f_t\) 定义的一致结构，与把 \(G\) 赋以拓扑 \(\mathfrak{G}\) 所得的拓扑群的左一致结构相同。

\begin{enumerate}
\def\labelenumi{\alph{enumi})}
\setcounter{enumi}{2}
\tightlist
\item
  在每个拓扑群 G 上，存在左不变伪度量组成的一个族，使得由该族定义的一致结构与 G 的左一致结构相同。
\end{enumerate}

\begin{enumerate}
\def\labelenumi{\arabic{enumi})}
\setcounter{enumi}{1}
\item
  设 \(\mathbf{G}\) 是一个左、右一致结构相重合的拓扑群。证明这个唯一的一致结构可以由一个同时左不变且右不变的伪度量族定义（利用第 III 章，§ 3，习题 3，证明：若 \(\mathbf{V}\) 是 \(e\) （\(\mathbf{G}\) 的单位元）的任一邻域，则 \(\mathbf{V}_0 = \bigcap_{x\in \mathbf{G}}x\mathbf{V}x^{-1}\) 是 \(e\) 的一个邻域）。
\item
  设 G 是一个拓扑群，\((f_{t})\) 是 G 上的定义 G 的左一致结构的一个左不变伪度量饱和族，且令 \(g_{t}(x)=f_{t}(e,x)\) 。设 H 是 G 的一个闭正规子群，对每个陪集 \(\dot{x}\in G/H\) ，令 \(h_{t}(\dot{x})=\inf_{x\in \dot{x}}g_{t}(x)\) ；证明：若 \(\bar{f}_{t}(\dot{x},\dot{y})=h_{t}(\dot{x}^{-1}\dot{y})\) ，则诸 \(\bar{f}_{t}\) 构成 G/H 上的一个左不变伪度量族，它定义 G/H 的左一致结构（按第 1 号的注记 2 的方式论证）。
\item
  证明赋值除环的拓扑是局部逆有界的（第 III 章，§ 6，习题 22）。
\end{enumerate}

¶ 5) 设 \(\omega\) 是有理数域 \(\mathbf{Q}\) 上的一个绝对值。a) 证明 \(\omega\) 由它在素数处的值唯一确定。b) 若存在一个素数 \(p\) 使 \(\omega(p) \leqslant 1\) ，证明 \(\omega(q) \leqslant 1\) 对每个其他素数 \(q\) 成立{[}给出 \(\omega(q^n)\) 的一个上界，办法是把 \(q^n\) 写成以 \(p\) 为基的尺度，然后令 \(n \to \infty\) {]}。

\begin{enumerate}
\def\labelenumi{\alph{enumi})}
\setcounter{enumi}{2}
\tightlist
\item
  若存在一个素数 p 使 \(\omega(p) < 1\) ，证明 \(\omega(q) = 1\) 对每个其他素数 q 成立{[}利用如下事实证明不可能有 \(\omega(q) < 1\) ：对每个整数 n \textgreater{} 0，存在两个有理整数 r 与 s 使得 \(1 = rp^{n} + sq^{n}\) {]} 。由此证明 \(\omega\) 此时是一个等价于 Q 上的 p 进绝对值的绝对值。d) 若对每个素数 p 有 \(\omega(p) > 1\) ，证明对任意两个素数 p 与 q，有
\end{enumerate}

\[
\frac {\log \omega (p)}{\log p} = \frac {\log \omega (q)}{\log q}
\]

{[}方法与 \(b\) 中相同{]}。由此证明在此情形 \(\omega(x) = |x|^{\rho}\) 对某个 \(\rho \leqslant 1\) 成立。

\begin{enumerate}
\def\labelenumi{\arabic{enumi})}
\setcounter{enumi}{5}
\item
  设 E 是除环 K 上的一个左向量空间，具有可数基 \((a_{n})\) 。设 \((r_{n})\) 是一个递减的、趋于 0 的数组成的序列 \textgreater0。对 E 的每个 \(x = \sum_{k} t_{k} a_{k} \neq 0\) ，令 \(\|x\| = r_{h}\) ，其中 h 是使 \(t_{k} \neq 0\) 的诸指标 k 中最小者；并令 \(\|0\| = 0\) 。证明 \(\|x - y\|\) 是 E 的加法群上的一个不变度量，且它在 E 上定义的拓扑与所选取的序列 \((r_{n})\) （递减且趋于 0）无关。由此证明：若把范数与赋范空间的定义（定义 5 与定义 6）推广到标量上的绝对值是平凡绝对值的情形，则命题 7 与定理 1 不再成立。
\item
  设 E 是赋值除环上的一个赋范空间。证明：若 E 中每个绝对可和族都在 E 中可和，则 E 是完备的。{[}设 \((x_{n})\) 是 E 中的一个柯西序列，考虑一个子序列 \((x_{n_{k}})\) （取自 \((x_{n})\) ），使得以 \(x_{n_{k+1}} - x_{n}\) 为通项的级数绝对收敛。{]}
\item
  给出 p 进数域 \(Q_{p}\) 中一个可和但不绝对可和的族的例子。{[}证明族 \((x_{t})_{t\in\mathbf{I}}\) （\(Q_{p}\) 的点组成的）是可和的，当且仅当 \(\lim x_{t}=0\) 关于 I 的有限子集的余集组成的滤子成立。{]}
\item
  映射 \(w\) （从环 \(\mathbf{A}\) 到 \(\mathbf{R}_+\) 中）称为 \(\mathbf{A}\) 上的半绝对值，如果它满足条件：(i) \(w(\mathrm{0}) = \mathrm{0}\) ；(ii) \(w(x - y) \leqslant w(x) + w(y)\) ；
\end{enumerate}

\begin{enumerate}
\def\labelenumi{(\roman{enumi})}
\setcounter{enumi}{2}
\tightlist
\item
  \(w(xy) \leqslant w(x)w(y)\) 。半绝对值 w 称为豪斯多夫的，如果 \(w(x) = 0\) 蕴含 x = 0。若 w 是 A 上的一个豪斯多夫半绝对值，则 \(w(x - y)\) 是 A 的加法群上的一个不变度量，从而定义一个与 A 的加法群结构相容的拓扑。把赋范代数的主要性质，特别是命题 13，推广到赋有半绝对值的环上。
\end{enumerate}

若 \(w\) 是 \(A\) 上一个非豪斯多夫的半绝对值，则集合 \(\overline{w}(0)\) 是一个双边理想 \(\mathfrak{a}\) （在 \(A\) 中）；若给 \(A\) 赋以由伪度量 \(w(x - y)\) 定义的拓扑，则该拓扑与 \(A\) 的环结构相容，且与 \(A\) 相伴的豪斯多夫空间是商环 \(A / \mathfrak{a}\) ，在它上面，函数 \(\overline{w}(\overline{x})\) ——对每个 \(\overline{x} \mod \mathfrak{a}\) ，它等于 \(w(x)\) （对所有 \(x \in \overline{x}\) ）的公共值——是一个称为与 \(w\) 相伴的豪斯多夫半绝对值；由 \(\overline{w}\) 定义的拓扑是商去 \(\mathfrak{a}\) 之后的 \(A\) 的拓扑。

¶ 10) a) 环 A 上的两个半绝对值 \(w_1, w_2\) 称为等价的，如果伪度量 \(w_1(x - y), w_2(x - y)\) 是等价的。证明：若 \(w\) 是 A 上的一个半绝对值，则 \(aw\) 与 \(w^{1/a}\) 是等价于 \(w\) 的半绝对值（对每个实数 \(a \geqslant 1\) ）。

\begin{enumerate}
\def\labelenumi{\alph{enumi})}
\setcounter{enumi}{1}
\tightlist
\item
  若 \(w_{i}\) （ \(1 \leqslant i \leqslant n\) ）是环 A 上的半绝对值，则函数 \(w = \sum_{i} w_{i}\) 与 \(w' = \sup_{i} w_{i}\) 是 A 上两个等价的半绝对值。若
\end{enumerate}

\[
\mathfrak {a} _ {i} = \overline {{{w}}} _ {i} ^ {1} (\mathrm{0}) \quad \text { and } \quad \mathfrak {a} = \overline {{{w}}} ^ {1} (\mathrm{0}),
\]

则有 \(\mathfrak{a} = \bigcap_{i} \mathfrak{a}_{i}\) 。若 \(A_i\) 表示商环 \(A/a_i\) （赋以与 \(w_i\) 相伴的豪斯多夫半绝对值）的完备化，证明 \(A/a_i\) （赋以与 \(w\) 相伴的豪斯多夫半绝对值）同构于乘积环 \(\prod_{i} A_i\) 的一个子环；证明（假设 \(A\) 有单位元 \(1\) ）这个环在 \(\prod_{i} A_i\) 中稠密，当且仅当诸 \(w_i\) 满足下列条件：对每个 \(\varepsilon > 0\) 及每个指标 \(i\) ，存在一个元素 \(x_i \in A\) ，使得 \(w_i(1 - x_i) \leqslant \varepsilon\) ，且 \(w_k(x_i) \leqslant \varepsilon\) 对每个指标 \(k \neq i\) 成立。

\begin{enumerate}
\def\labelenumi{\arabic{enumi})}
\setcounter{enumi}{10}
\tightlist
\item
  设 A 是赋有豪斯多夫半绝对值的一个环，并设 A 有单位元，且关于由该半绝对值定义的拓扑是完备的。证明 A 的每个极大理想都是闭的（利用命题 13）。
\end{enumerate}

¶ 12) 设 K 是一个非离散的豪斯多夫拓扑除环，\(\varphi: \mathbf{K} \to \mathbf{R}_{+}\) 是一个映射，满足 \(\varphi(0) = 0\) ，\(\varphi(xy) = \varphi(x)\varphi(y)\) ，且使得：若 \(V_{n}\) 表示所有 \(x \in K\) 中满足 \(\varphi(x) \leqslant 1/n\) 者组成的集合，则诸集 \(V_{n}\) 构成 K 中 \(0\) 的一个基本邻域系。

\begin{enumerate}
\def\labelenumi{\alph{enumi})}
\tightlist
\item
  证明存在一个实数 \(a > 0\) ，使得对所有 \(x \in \mathbf{K}\) ，
\end{enumerate}

\[
\varphi (1 + x) \leqslant a (1 + \varphi (x))
\]

{[}否则，将存在 K 的点组成的序列 \((x_{n})\) ，使得 \((1 + x_n)^{-1}\) 与 \(x_{n}(1 + x_{n})^{-1}\) 都趋于零{]}。由此证明：若 \(\psi (x) = (\varphi (x))^{\alpha}\) ，则有

\[
\psi (x + y) \leqslant 2 \sup (\psi (x), \psi (y))
\]

对充分小的 \(\alpha\) 成立。

b) 证明若 \(n = 2^p\) ，则有 \(\psi\left(\sum_{i=1}^{n} x_i\right) \leqslant n \sup (\psi(x_i))\) ，并推出对每个整数 \(m > 0\) 有 \(\psi(m) \leqslant 2m\) 。

\begin{enumerate}
\def\labelenumi{\alph{enumi})}
\setcounter{enumi}{2}
\tightlist
\item
  由 b) 推出：对每个 \(n = 2^p\) 及每个 \(x \in \mathbf{K}\) ，有 \(\psi((1 + x)^{n-1}) \leqslant 2n(1 + \psi(x))^{n-1}\) ，从而 \(\psi\) 是 \(\mathbf{K}\) 上的一个定义 \(\mathbf{K}\) 的拓扑的绝对值。
\end{enumerate}

\begin{enumerate}
\def\labelenumi{\arabic{enumi})}
\setcounter{enumi}{12}
\tightlist
\item
  设 K 是一个非离散的豪斯多夫拓扑除环。设 R 是所有 \(x \in K\) 中满足 \(\lim_{n \to \infty} x^n = 0\) 者组成的集合，N 是集合 \(R \cup R^{-1}\) 的余集。证明：K 上存在一个定义 K 的拓扑的绝对值，当且仅当 (i) R 在 K 中是开的；(ii) 对 K 中 0 的每个邻域 V，存在 K 中 0 的一个邻域 U，使得 \(R.U \subset V\) ；且 (iii) 若 \(x \in R\) 且 \(y \in R \cup N\) ，则 \(yx \in R\) 。
\end{enumerate}

为证明这些条件是充分的，依次证明：a) N 是 K 的非零元素组成的乘法群 \(K^{*}\) 的一个正规子群。

\begin{enumerate}
\def\labelenumi{\alph{enumi})}
\setcounter{enumi}{1}
\item
  在商群 \(K^{*}/N\) 中，规定 \(\dot{x} \leqslant \dot{y}\) ，如果存在 \(x \in \dot{x}\) 与 \(y \in \dot{y}\) 使得 \(yx^{-1} \in R \cup N\) ；证明这个关系是一个与 \(K^{*}/N\) 的群结构相容的序关系，且如此定义的有序群同构于加法群 R 的一个子群（利用第 V 章，§ 3，习题 1）。
\item
  借助习题 12 完成证明。
\end{enumerate}

特别地，非离散局部紧拓扑域的拓扑可以由一个绝对值定义。

